% --------------------------------------------------------------------------- %
% Poster for the ECCS 2011 Conference about Elementary Dynamic Networks.      %
% --------------------------------------------------------------------------- %
% Created with Brian Amberg's LaTeX Poster Template. Please refer for the     %
% attached README.md file for the details how to compile with `pdflatex`.     %
% --------------------------------------------------------------------------- %
% $LastChangedDate:: 2011-09-11 10:57:12 +0200 (V, 11 szept. 2011)          $ %
% $LastChangedRevision:: 128                                                $ %
% $LastChangedBy:: rlegendi                                                 $ %
% $Id:: poster.tex 128 2011-09-11 08:57:12Z rlegendi                        $ %
% --------------------------------------------------------------------------- %
\documentclass[a0paper,portrait]{baposter}
\newcommand{\irptitle}{Remaining Anonymous when using the Bitcoin Protocol}
\usepackage[style=ieee]{biblatex}
\addbibresource{i-d.bib}
\addbibresource{report.bib}
\addbibresource{rfc.bib}

\newcommand\irptopic{Bitcoin}
\newcommand\irpauthor{Thomas A. Grainger}

\newcommand\MYhyperrefoptions{bookmarks=true,bookmarksnumbered=true,
pdfpagemode={UseOutlines},plainpages=false,pdfpagelabels=true,
hidelinks,
pdftitle={\irptitle},%<!CHANGE!
pdfsubject={\irptopic},%<!CHANGE!
pdfauthor={Thomas A. Grainger},%<!CHANGE!
pdfkeywords={Bitcoin, Anonymity, Cryptography, crypto-currencies, Peer to Peer}}
\usepackage{setspace}

\renewcommand{\biblabelsep}{2em}
\renewcommand{\bibitemsep}{0.2ex}
\renewcommand*{\bibfont}{\tiny}


\usepackage{graphicx}
\usepackage{multicol}
\setlength{\columnsep}{1.5em}
\setlength{\columnseprule}{0mm}

\usepackage{enumitem}
\setdescription{leftmargin=*}
%\newcommand\MYhyperrefoptions{bookmarks=true,bookmarksnumbered=true,
%pdfpagemode={UseOutlines},plainpages=false,pdfpagelabels=true,
%hidelinks,
%pdftitle={\irptitle},%<!CHANGE!
%pdfsubject={Bitcoin},%<!CHANGE!
%pdfauthor={Thomas A. Grainger},%<!CHANGE!
%pdfkeywords={Bitcoin, Anonymity, Cryptography, crypto-currencies, Peer to Peer}}

%\usepackage[\MYhyperrefoptions,pdftex]{hyperref}

\usepackage{relsize}    % For \smaller
\hyphenation{op-tical net-works semi-conduc-tor}
\usepackage{amssymb}
\usepackage{amsmath}
%%% Global Settings %%%%%%%%%%%%%%%%%%%%%%%%%%%%%%%%%%%%%%%%%%%%%%%%%%%%%%%%%%%

%%% Color Definitions %%%%%%%%%%%%%%%%%%%%%%%%%%%%%%%%%%%%%%%%%%%%%%%%%%%%%%%%%

\definecolor{bordercol}{RGB}{40,40,40}
\definecolor{headercol1}{RGB}{186,215,230}
\definecolor{headercol2}{RGB}{80,80,80}
\definecolor{headerfontcol}{RGB}{0,0,0}
\definecolor{boxcolor}{RGB}{186,215,230}

%%%%%%%%%%%%%%%%%%%%%%%%%%%%%%%%%%%%%%%%%%%%%%%%%%%%%%%%%%%%%%%%%%%%%%%%%%%%%%%%
%%% Utility functions %%%%%%%%%%%%%%%%%%%%%%%%%%%%%%%%%%%%%%%%%%%%%%%%%%%%%%%%%%

%%% Save space in lists. Use this after the opening of the list %%%%%%%%%%%%%%%%
\newcommand{\compresslist}{
    \setlength{\itemsep}{1pt}
    \setlength{\parskip}{0pt}
    \setlength{\parsep}{0pt}
}

%%%%%%%%%%%%%%%%%%%%%%%%%%%%%%%%%%%%%%%%%%%%%%%%%%%%%%%%%%%%%%%%%%%%%%%%%%%%%%%
%%% Document Start %%%%%%%%%%%%%%%%%%%%%%%%%%%%%%%%%%%%%%%%%%%%%%%%%%%%%%%%%%%%
%%%%%%%%%%%%%%%%%%%%%%%%%%%%%%%%%%%%%%%%%%%%%%%%%%%%%%%%%%%%%%%%%%%%%%%%%%%%%%%

\begin{document}

%%% Setting Background Image %%%%%%%%%%%%%%%%%%%%%%%%%%%%%%%%%%%%%%%%%%%%%%%%%%
\background{
    \begin{tikzpicture}[remember picture,overlay]%
    \draw (current page.north west)+(-2em,2em) node[anchor=north west]
    {\includegraphics[height=1.1\textheight]{pix/background}};
    \end{tikzpicture}
}

%%% General Poster Settings %%%%%%%%%%%%%%%%%%%%%%%%%%%%%%%%%%%%%%%%%%%%%%%%%%%
%%%%%% Eye Catcher, Title, Authors and University Images %%%%%%%%%%%%%%%%%%%%%%
\begin{poster}{
    grid=false,
    % Option is left on true though the eyecatcher is not used. The reason is
    % that we have a bit nicer looking title and author formatting in the headercol
    % this way
    eyecatcher=true, 
    borderColor=bordercol,
    headerColorOne=headercol1,
    headerColorTwo=headercol2,
    headerFontColor=headerfontcol,
    % Only simple background color used, no shading, so boxColorTwo isn't necessary
    boxColorOne=boxcolor,
    headershape=roundedright,
    headerfont=\Large\sf\bf,
    textborder=rectangle,
    background=user,
    headerborder=open,
    boxshade=plain
}
%%% Eye Cacther %%%%%%%%%%%%%%%%%%%%%%%%%%%%%%%%%%%%%%%%%%%%%%%%%%%%%%%%%%%%%%%
{
% The logos are compressed a bit into a simple box to make them smaller on the result
% (Wasn't able to find any bigger of them.)
    \includegraphics[height=3em]{img/ecs_descriptor}

}
%%% Title %%%%%%%%%%%%%%%%%%%%%%%%%%%%%%%%%%%%%%%%%%%%%%%%%%%%%%%%%%%%%%%%%%%%%
{\sf\bf
    \irptitle
}
%%% Authors %%%%%%%%%%%%%%%%%%%%%%%%%%%%%%%%%%%%%%%%%%%%%%%%%%%%%%%%%%%%%%%%%%%
{
    \vspace{1em} Thomas Grainger\\
    {\smaller Email: t.grainger@ecs.soton.ac.uk}
}
%%% Logo %%%%%%%%%%%%%%%%%%%%%%%%%%%%%%%%%%%%%%%%%%%%%%%%%%%%%%%%%%%%%%%%%%%%%%
{
% The logos are compressed a bit into a simple box to make them smaller on the result
% (Wasn't able to find any bigger of them.)
    \includegraphics[height=3em]{img/University_of_Southampton}
}

\headerbox{Introduction}{name=introduction,column=0,row=0}{
While cryptographic digital currencies have previously been discussed as mathematical curiosities, Bitcoin\cite{satoshi}, a peer-to-peer electronic currency system has become the focus of significant interest and research effort. Investigated are methods used for de-anonymizing or revealing the identity of users of the Bitcoin network and how these methods can be contravened as well as proposing a novel automated solution to the issue of de-anonymizing techniques.
}

\headerbox{Existing Systems}{name=existing,column=0,below=introduction}{

\begin{description}
    \item[Central Banking Authorities] These fully central systems, such as PayPal, or SEPA handle transactions by debiting one user account and crediting another~\cite{paypal}.
    \item[Unique Token (NetCash)] To create NetCash~\cite{netcash} currency, the user and central authority agree on a unique token in exchange for some currency or resource. To transfer value a user sends a token to another user. Receiving users send the token to the exchange, to be swapped for a new spendable token invalidating the original token.
    \item[Blind-signed Chaumian e-Cash] Uses blind signatures. A user, Alice, creates a token $x$ and blinds it using function $c$.  A central authority will sign this blinded message using signing function $s$ in exchange for a fee, creating $s(c(x))$.  The authority then returns the blind signed message to Alice. Alice then removes the blind, leaving the signature $s(x)$ using $c^{-1}(s(c(x)))$\cite{chaum}.
    \item[B-Money digital signature chain]B-Money proposed by~\textcite{b-money}. Transaction messages are created using public key cryptography and the chain of digital signatures can be traced back to the original creation of value. Money is not sent per-se but re-assigned. To prevent double spending all users maintain a live copy of all transactions using a theoretical ``synchronous and unjammable anonymous broadcast channel''.
\end{description}
    \includegraphics[width=\textwidth]{img/Bitcoin_Transaction_Visual}
}

\headerbox{Acknowledgements}{name=acknowledgements,column=0,below=existing, above=bottom}{

\vspace{-0.4em}     % Save some space at the beginning
The author would like to thank Dr~Tim~Chown for his
support and supervision throughout this project, the mysterious Satoshi Nakamoto for the creation of the Bitcoin distributed currency system, and Gregory (gmaxwell) Maxwell and the other users of the freenode
\#bitcoin-dev Internet Relay Chat (IRC) channel.
}



\headerbox{The Bitcoin Protocol}{name=bitcoin,span=2,column=1,row=0}{
\begin{multicols}{2}
The Bitcoin protocol and client software provides a concrete implementation of an altered version of the original B-Money proposal.  The only significant novel deviation is that \textcite{satoshi} has substituted the ``synchronous and unjammable anonymous broadcast channel''~\cite{b-money} with a decentralized time stamp server.

The Bitcoin peer network, create a Bitcoin block chain: a chain of blocks of transactions, each referring to the previous by hash value. Each block has a ``proof-of-work'' requirement: the double SHA-256 hash value of each block must not exceed a value determined by the current `difficulty'. This means to replace a block in the chain all later blocks must also be recreated.
\end{multicols}
\vspace{-1.5em}
\begin{center}
    \includegraphics[width=0.40\linewidth]{img/dyn/price}
    \includegraphics[width=0.40\linewidth]{img/dyn/speed-lin-ever}
\end{center}
}

\headerbox{De-Anonymization Techniques}
{name=deanon,span=2,column=1,below=bitcoin}{
\begin{multicols}{2}
In Bitcoin all transactions must be broadcast publicly, however, \textcite{satoshi} maintains that ``privacy can still be maintained by breaking the flow of information [...] by keeping public keys anonymous''.

\begin{description}
\item[Multi-input Transactions]
If a user cannot fulfill the value of a transaction by spending a single previous output, the client will combine multiple inputs, where the user has the ability to fulfill their spend conditions.  ``It is therefore straightforward to conclude that if these (inputs) are (spendable) by different addresses, then the input addresses belong to the same user''\cite{reid-anon}.

\item[Change addresses]
Because each output must be redeemed in it's entirety transactions of values smaller than the smallest output available to a wallet, the Bitcoin client uses ``change addresses'' so as to retain the remaining value.\textcite{eval-priv} propose a heuristic: if money is sent to two or more addresses and only one has not been used before, that address is the change addresses controlled by the user.
\end{description}
By combining the knowledge of how keys are grouped into wallets with knowledge from outside the protocol the identity of the user behind the wallet can be discovered\cite{reid-anon}.
\includegraphics[width=\columnwidth]{img/group_addresses}
\end{multicols}
}

\headerbox{Anonymization Techniques}
{name=anon,span=2,column=1,below=deanon}{
\begin{multicols}{2}
Various methods exist to avoid de-anonymization:
\begin{description}
\item[Fully Trusted Mixing]
A mixing service accepts Bitcoin transactions and automatically sends a new transaction to a newly generated Bitcoin address.  These services often charge fees and require a significant amount of trust.
\item[ZeroCoin]
\textcite{zerocoin} propose a method of avoiding de-anonymization without a central trusted party, by defining new transaction types that take advantage of a zero knowledge proof (ZKP) that only proves knowledge of: a Zerocoin in the block chain. The actual serial number used to generate it. This is used to spend money without revealing which serial number was added to the block chain.

\item[In-Protocol Mixing]
Because the conditions of a set of previous outputs can be fulfilled partially, multiple users can cooperate to create a single transaction that causes their public keys to appear linked in a single wallet to taint analysis, without revealing their private keys to the group.  As such, the claim made by \textcite{satoshi}, and key assumption made by \textcite{reid-anon}, that some linking in multi-input transactions is unavoidable which necessarily reveals that inputs were owned by the same user, is false.

\end{description}
\end{multicols}
}

\headerbox{References}{name=references,span=2,column=1,below=anon, above=bottom}{
\vspace{-0.4em}                     % Save some space at the beginning
\renewcommand{\section}[2]{\vskip 0.05em}   % Omit "References" title
\printbibliography
}

\end{poster}
\end{document}
